\section*{Repository and code availability}
The full source code is maintained in a public Git repository to make the
workflow reproducible and auditable.

Repository URL: \url{https://github.com/Li-Wen-Di/BaikovGF}

The repository includes the Wolfram Language package, the external backend,
and example notebooks and tests. A clone of the repository is sufficient to
run the package on a new machine, subject to the software dependencies listed
below.

\section*{Implementation language (Mathematica/Wolfram Language)}
The implementation is written in Wolfram Language (Mathematica), with the main
package file stored under \texttt{BaikovGF/BaikovGF.wl}. The code is distributed
as plain text \texttt{.wl} and \texttt{.mma} files so it can be inspected and
versioned in Git without extra tooling.

To load the package from a local clone:
\begin{verbatim}
projectRoot = "/path/to/this/repository";
Get[FileNameJoin[{projectRoot, "BaikovGF", "BaikovGF.wl"}]];
Needs["BaikovGF`"];
\end{verbatim}

This loads the symbols into the current session and makes the public interface
available for analytic workflows.

\section*{External dependency: python-flint}
Some simplification stages are delegated to a lightweight Python backend that
uses \texttt{python-flint} for fast rational operations. The package prefers a
project-local virtual environment named \texttt{.venv-flint} in the repository
root. If it is not present, the package falls back to the system Python and
checks whether the \texttt{flint} module can be imported.

We provide platform-specific setup scripts that create \texttt{.venv-flint} and
install \texttt{python-flint}:
\begin{itemize}
  \item Windows: \texttt{external/setup\_flint\_env.ps1}
  \item Linux: \texttt{external/setup\_flint\_env\_linux.sh}
  \item macOS: \texttt{external/setup\_flint\_env\_macos.sh}
\end{itemize}

These scripts are designed to be run from the repository root and do not require
manual path edits. They install a private dependency set within the repository
so that the environment can be recreated by other users without changing system
packages.

\section*{Backend availability check}
The package provides a built-in check that verifies the external backend. It
detects the Python executable, tests \texttt{import flint}, and runs a minimal
round-trip simplification:
\begin{verbatim}
CheckBaikovExternalBackend[]
\end{verbatim}

If the check reports an error, the most common fix is to run the appropriate
setup script listed above and re-run the check.

\section*{Public interface (core functions)}
The following functions constitute the public interface used in the paper:
\begin{itemize}
  \item \texttt{SP[p, q]}: symbolic scalar product used to define kinematics.
  \item \texttt{CreateFeynmanIntegral[...]}: construct a family association from
    loop/external momenta, propagators, and kinematic rules.
  \item \texttt{BuildBaikovGeneratingFunction[...]}: build a generating-function
    object for the requested target powers.
  \item \texttt{ExtractBaikovCoefficient[...]}: extract a single unsimplified
    coefficient from a generating function.
  \item \texttt{ExtractBaikovCoefficientBatch[...]}: extract a list of
    coefficients in a single call.
  \item \texttt{SimplifyBaikovCoefficient[...]}: simplify one coefficient, using
    internal algebra and (optionally) the external backend.
  \item \texttt{SimplifyBaikovCoefficientBatch[...]}: simplify a list of
    coefficients.
  \item \texttt{CheckBaikovExternalBackend[]}: verify that the Python/FLINT
    backend is available and functional.
\end{itemize}

\section*{Minimal reproducible workflow}
The following outline shows the minimal sequence used in analytic runs:
\begin{verbatim}
projectRoot = "/path/to/this/repository";
Get[FileNameJoin[{projectRoot, "BaikovGF", "BaikovGF.wl"}]];
Needs["BaikovGF`"];

family = CreateFeynmanIntegral[
  "LoopMomenta" -> {...},
  "ExternalMomenta" -> {...},
  "Propagators" -> {...},
  "AuxiliaryPropagators" -> {...},
  "KinematicRules" -> {...},
  "DimensionSymbol" -> d
];

gf = BuildBaikovGeneratingFunction[family, targetPowers];
raw = ExtractBaikovCoefficient[gf, sourcePowers];
res = SimplifyBaikovCoefficient[raw];

res["Coefficient"]
\end{verbatim}

This sequence isolates the construction, extraction, and simplification stages
that are reported in the paper. When the external backend is available, the
final rational simplification is delegated automatically.
